\documentclass[11pt,a4paper]{article}

\usepackage[utf8]{inputenc}
\usepackage[T1]{fontenc}
\usepackage[spanish,es-tabla]{babel}
\usepackage{amsmath,amssymb,amsthm,mathtools}
\usepackage{booktabs}
\usepackage{enumitem}
\usepackage{geometry}
\usepackage{setspace}
\usepackage{hyperref}
\usepackage{microtype}

\geometry{margin=1in}
\onehalfspacing

\hypersetup{
  pdftitle={Nota formal: Esbozo categorico de conjunciones ordinales},
  pdfauthor={Johel Padilla-Villanueva}
}

\theoremstyle{definition}
\newtheorem{definition}{Definici\'on}[section]
\newtheorem{remark}{Observaci\'on}[section]
\theoremstyle{plain}
\newtheorem{proposition}{Proposici\'on}[section]
\newtheorem{theorem}{Teorema}[section]
\newtheorem{conjecture}{Conjetura}[section]
\theoremstyle{remark}
\newtheorem{scholium}{Scholium}[section]

\newcommand{\Res}{\mathrm{Res}}
\newcommand{\Rpair}{\mathcal{R}_{\mathrm{pair}}}
\newcommand{\Rind}{\mathcal{R}_{\mathrm{ind}}}
\newcommand{\Syn}{\mathrm{Syn}}
\newcommand{\KL}{D_{\mathrm{KL}}}

\title{\textbf{Nota formal}\\[0.4em]
\large Esbozo categ\'orico y algebraico\\
de los niveles de conjunci\'on ordinal\\[0.6em]
\normalsize (Punto 6 del roadmap --- opcional; asume Notas 1--2)}

\author{
Johel Padilla-Villanueva\\
\textit{Department of Environmental Health, University of Puerto Rico}\\
\texttt{johelpadilla@gmail.com}
}

\date{2026-08-01 \quad$\cdot$\quad v0.1}

\begin{document}
\maketitle

\begin{abstract}
Esbozo \emph{elegante y m\'inimo} ---no un tratado de categor\'ias---
para expresar los niveles de conjunci\'on como funtores y el Nivel~3
como obstrucci\'on a factorizar por un funtor de olvido hacia estructuras
de pares. El objetivo del punto~6 es elevar la claridad estructural del
paradigma, no sustituir las Notas 1--5 ni introducir metaf\'isica.
\end{abstract}

\tableofcontents
\vspace{1em}

\section{Advertencia de oficio}
\label{sec:warn}

\begin{scholium}
Este punto es \textbf{opcional}. Si alguna construcci\'on choca con las
definiciones de las Notas 1--2, mandan las Notas 1--2.
No se usa lenguaje categ\'orico para simular profundidad ontológica
(M-I-1).
\end{scholium}

\section{Categor\'ia de procesos de s\'imbolos}
\label{sec:cat}

\begin{definition}[Categor\'ia $\mathbf{Sym}$]
\label{def:Sym}
\begin{itemize}
  \item \textbf{Objetos:} procesos (o leyes) $P$ sobre $\Sigma^d$ con
  $d\ge 1$ y $\Sigma=S_m$ fijo en un contexto, o pares
  $(P,d,m)$.
  \item \textbf{Morfismos:} aplicaciones que preservan la estructura
  ordinal relevante. Ejemplos:
  \begin{enumerate}[label=(\alph*)]
    \item reetiquetado de canales $\sigma\in S_d$;
    \item transformaciones mon\'otonas estrictas por canal (inducen
    $\mathrm{id}$ en s\'imbolos);
    \item marginaciones $P\mapsto P_A$ a $A\subset\{1,\ldots,d\}$.
  \end{enumerate}
\end{itemize}
Una elecci\'on austera: morfismos $=$ permutaciones de canales +
isometr\'ias mon\'otonas (entonces muchos morfismos son isomorfismos
sobre $S$).
\end{definition}

\begin{definition}[Funtor de s\'imbolos]
\label{def:BP}
$\mathrm{BP}:\mathbf{Series}\to\mathbf{Sym}$ env\'ia trayectorias
reales multivariantes a su proceso Bandt--Pompe. Es el puente
``datos $\to$ objetos ordinales''.
\end{definition}

\section{Estructuras de profundidad como funtores}
\label{sec:depth}

\begin{definition}[Funtor de pares]
\label{def:Pi2}
$\Pi_2:\mathbf{Sym}\to\mathbf{Pair}$ env\'ia $P$ a la familia de
bipolares $\{P_{ij}\}_{i<j}$ (o a la I-proyecci\'on $P^{(2)}\in\Rpair$).
$\mathbf{Pair}$ es la categor\'ia de colecciones de leyes bipolares
consistentes (o de modelos de pares).
\end{definition}

\begin{definition}[Funtor de independencia]
\label{def:Pi1}
$\Pi_{\mathrm{ind}}:\mathbf{Sym}\to\mathbf{Marg}$ env\'ia $P$ a
$(P_1,\ldots,P_d)$ o a $P_{\mathrm{ind}}=\prod P_i$.
\end{definition}

\begin{proposition}[Factorizaci\'on de olvido]
\label{prop:forget}
Hay un tri\'angulo natural
\begin{equation}
\mathbf{Sym}
\xrightarrow{\Pi_2}
\mathbf{Pair}
\xrightarrow{\Pi_{\mathrm{ind}}'}
\mathbf{Marg},
\end{equation}
y $\Pi_{\mathrm{ind}}$ factoriza (hasta isomorfismo natural) por
$\Pi_2$ en el sentido de que la independencia total olvida m\'as que el
olvido a pares.
\end{proposition}

\begin{definition}[Claims como funtores valuados]
\label{def:claims_fun}
\begin{align}
\mathcal{C}_1:\mathbf{Sym}&\to\mathbf{Meas},&
&P\mapsto\Phi_1(P)\ \text{(o proceso $\Phi_1(t)$)},\\
\mathcal{C}_2:\mathbf{Sym}&\to\mathbf{Meas},&
&P\mapsto\Phi_2(P),\\
\mathcal{C}_3:\mathbf{Sym}&\to\mathbf{Meas},&
&P\mapsto\Res_{\mathrm{pair}}(P).
\end{align}
$\mathbf{Meas}$ es la categor\'ia de espacios de medida / $[0,\infty)$-valores
con morfismos mon\'otonos adecuados.
\end{definition}

\begin{theorem}[L3 como obstrucci\'on al olvido]
\label{thm:obstruction}
$\Res_{\mathrm{pair}}(P)=0$ sii $P$ est\'a en la imagen esencial de una
secci\'on de $\Pi_2$ (i.e., $P$ es un modelo de pares). Por tanto
$\mathcal{C}_3(P)$ mide la \textbf{obstrucci\'on} a que $P$ descienda a
$\mathbf{Pair}$ sin p\'erdida (en divergencia KL).
\end{theorem}

\begin{proof}
Por definici\'on de $\Res_{\mathrm{pair}}=\inf_{Q\in\Rpair}\KL(P\|Q)$ y
existencia de $P^{(2)}$: se anula sii $P=P^{(2)}\in\Rpair$.
\end{proof}

\begin{scholium}
Esta es la formulaci\'on categ\'orica de ``sinergia irreducible'': no es
un n\'umero m\'agico, es la obstrucci\'on al funtor de olvido $\Pi_2$.
\end{scholium}

\section{Poset de profundidad}
\label{sec:poset}

\begin{definition}[Poset de claims]
\label{def:poset}
Sea $\mathbf{D}=\{1,2,3\}$ con orden $1\prec 2\prec 3$ interpretado como
\emph{exigencia estructural} (anidamiento Tipo~II, Nota~1), no como
implicaci\'on de indicadores.
Un \textbf{estado de profundidad} es un elemento de $\mathbf{D}$ o una
distribuci\'on $(f_1,f_2,f_3)$ en el simplejo (share del reloj, Nota~3).
\end{definition}

\begin{definition}[Motor de pesos como funtor de r\'egimen]
\label{def:weight_fun}
Sea $\mathbf{Reg}$ la categor\'ia de cues $\lambda$ (objetos
$\mathbb{R}_+$, morfismos mon\'otonos que preservan orden de r\'egimen).
Un motor admisible es un funtor (o simplemente un mapa mon\'otono)
\begin{equation}
\alpha:\mathbf{Reg}\to\mathbb{R}_{+}^{3}
\end{equation}
con la monotom\'ia de profundidad de la Nota~3. El RECD es la
composici\'on
\begin{equation}
(\mathcal{C}_1,\mathcal{C}_2,\mathcal{C}_3)\otimes\alpha(\lambda)
\ \text{en el sentido de producto interno}\ 
\sum_\ell\alpha_\ell\mathcal{C}_\ell.
\end{equation}
\end{definition}

\section{\'Algebra de interacciones}
\label{sec:alg}

\begin{remark}[Familia exponencial]
En la familia log-lineal
\begin{equation}
\log P(s)
=\sum_i\theta_i(s_i)
+\sum_{i<j}\theta_{ij}(s_i,s_j)
+\sum_{i<j<k}\theta_{ijk}(s_i,s_j,s_k)
+\cdots,
\end{equation}
$\Rpair$ es el submodelo $\theta_{ijk\ldots}=0$. Los coeficientes de
orden $\ge 3$ son el contenido algebraico de L3; $\Res_{\mathrm{pair}}$
es una norma (KL) del residuo al proyectar esos coeficientes a cero.
\end{remark}

\begin{proposition}[Lectura algebraica de excess3]
\label{prop:alg_excess}
$\Syn$ es un \emph{proxy num\'erico} del tama\~no de interacciones de
orden $\ge 3$ tras una resta heur\'istica de pares; no es el vector
$(\theta_{ijk})$. PID intenta descomponer contribuciones de informaci\'on
en un lattice; el esbozo categ\'orico de esta nota no reemplaza PID.
\end{proposition}

\section{Diagrama resumen}
\label{sec:diag}

\begin{center}
\begin{tabular}{@{}c@{}}
$\mathbf{Series}$
\ $\xrightarrow{\mathrm{BP}}$\
$\mathbf{Sym}$
\ $\xrightarrow{\Pi_2}$\
$\mathbf{Pair}$
\ $\xrightarrow{\Pi_{\mathrm{ind}}'}$\
$\mathbf{Marg}$\\[0.6em]
\hspace{3.2cm}
$\Big\downarrow \mathcal{C}_1,\mathcal{C}_2$
\hspace{0.8cm}
$\Big\downarrow 0$
\hspace{1.2cm}
$\Big\downarrow 0$\\[0.3em]
\hspace{3.2cm}
$\mathbf{Meas}$
\hspace{0.3cm}
(obstrucci\'on $\mathcal{C}_3=\Res_{\mathrm{pair}}$)
\end{tabular}
\end{center}

\section{Canon a\~nadido (C35--C38)}
\label{sec:canon}

\begin{enumerate}[label=\textbf{C\arabic*.}]
  \setcounter{enumi}{34}
  \item $\Pi_2=$ funtor de olvido a pares; L3 $=$ obstrucci\'on a
  $\Pi_2$ (Thm.~\ref{thm:obstruction}).
  \item Claims $\mathcal{C}_\ell$ como funtores $\mathbf{Sym}\to\mathbf{Meas}$.
  \item Anidamiento Tipo~II $=$ orden en el poset $\mathbf{D}=\{1\prec 2\prec 3\}$.
  \item El lenguaje categ\'orico es opcional y subordinado a Notas 1--2.
\end{enumerate}

\vspace{1.5em}
\noindent\rule{\textwidth}{0.4pt}\\[0.5em]
\begin{center}
\em Fin Nota formal punto 6 v0.1 (esbozo).\\
Con esto se completa el arco 1--6 del roadmap te\'orico en versi\'on v0.1.
\end{center}

\begin{thebibliography}{99}
\bibitem{Padilla2026note1}
J.~Padilla-Villanueva. Notas formales 1--5. Fundaci\'on RECD, 2026.
\bibitem{MacLane}
S.~Mac Lane. \emph{Categories for the Working Mathematician}. Springer.
\end{thebibliography}

\end{document}
